\documentclass[11pt,letterpaper]{article}

\usepackage[utf8]{inputenc}
\usepackage[spanish]{babel}
\usepackage[margin=2.5cm]{geometry}
\usepackage{amsmath,amssymb,amsthm}
\usepackage{enumitem}
\usepackage{titlesec}
\usepackage{array}
\usepackage{hyperref}
\usepackage{xcolor}
\usepackage{listings}
\usepackage{tikz}
\usetikzlibrary{arrows.meta,positioning,calc}

\definecolor{unicaucaverde}{RGB}{0,90,60}
\hypersetup{colorlinks=true,linkcolor=unicaucaverde,urlcolor=unicaucaverde}

\titleformat{\section}{\normalfont\Large\bfseries\color{unicaucaverde}}{\thesection}{1em}{}
\titleformat{\subsection}{\normalfont\large\bfseries}{\thesubsection}{1em}{}

\theoremstyle{definition}
\newtheorem{definicion}{Definición}[section]
\newtheorem{ejemplo}[definicion]{Ejemplo}
\newtheorem{propiedad}[definicion]{Propiedad}

\lstdefinestyle{cstyle}{
  language=C,
  basicstyle=\ttfamily\small,
  keywordstyle=\color{unicaucaverde}\bfseries,
  commentstyle=\color{gray}\itshape,
  stringstyle=\color{orange!70!black},
  numbers=left,
  numberstyle=\tiny\color{gray},
  frame=single,
  breaklines=true,
  showstringspaces=false,
  tabsize=2,
  captionpos=b
}
\lstdefinestyle{vstyle}{
  language=Verilog,
  basicstyle=\ttfamily\small,
  keywordstyle=\color{unicaucaverde}\bfseries,
  commentstyle=\color{gray}\itshape,
  stringstyle=\color{orange!70!black},
  numbers=left,
  numberstyle=\tiny\color{gray},
  frame=single,
  breaklines=true,
  showstringspaces=false,
  tabsize=2,
  captionpos=b
}
\lstset{style=cstyle}

\title{\textcolor{unicaucaverde}{\Large Notas de clase --- Semana 1}\\[4pt]
  \large Sistemas embebidos en FPGA: del sistema Nios II en Quartus al software en Nios II EDS}
\author{Énfasis --- Universidad del Cauca}
\date{2026}

\begin{document}
\maketitle
\tableofcontents

\section{Procesadores soft-core vs. hard-core; arquitectura Nios II}

\begin{definicion}[Soft processor]
Un \emph{soft processor} (o \emph{soft core}) es un procesador completo --- ruta de datos, unidad
de control, banco de registros --- descrito en HDL y sintetizado dentro del tejido lógico
reconfigurable de una FPGA, en lugar de venir fabricado como bloque fijo en el silicio
(\emph{hard core}, como los núcleos ARM de un SoC Zynq o Cyclone V SoC). Nios II, de
Intel/Altera, es el soft processor de referencia para las familias Cyclone.
\end{definicion}

\begin{propiedad}[Por qué usar un soft core en vez de un microcontrolador dedicado]
Un microcontrolador dedicado suele ser más barato y de menor consumo para una tarea fija. Un soft
core tiene sentido cuando el procesador necesita convivir, dentro del mismo chip, con lógica digital
hecha a la medida --- controladores propios, aceleradores en hardware, interfaces no estándar ---
comunicados por el mismo bus interno: el diseñador decide en qué parte del sistema resolver un
problema en software (más flexible, más lento) y en qué parte resolverlo en lógica dedicada (más
rápido, fijo), y puede mover esa frontera sin cambiar de chip. El costo es área lógica de la FPGA
(el soft core ocupa LEs que ya no están disponibles para lógica propia) y una frecuencia de reloj
bastante menor que la de un hard core equivalente.
\end{propiedad}

\begin{definicion}[Núcleos Nios II]
Nios II es una arquitectura RISC de 32 bits, configurable en tres variantes que intercambian área
por desempeño:
\begin{itemize}[leftmargin=1.2cm]
  \item \textbf{Nios II/e (Economy)}: la más pequeña; sin \emph{caches}, sin predicción de saltos,
        ejecuta cada instrucción en varios ciclos. Útil cuando el área de la FPGA es el recurso más
        escaso.
  \item \textbf{Nios II/s (Standard)}: punto intermedio; pipeline de 5 etapas, sin caches de
        datos/instrucciones.
  \item \textbf{Nios II/f (Fast)}: la de mejor desempeño; caches de instrucciones y datos,
        predicción dinámica de saltos, multiplicador/divisor en hardware.
\end{itemize}
Para un primer sistema mínimo --- el de esta semana --- basta el núcleo \texttt{/e}: el objetivo es
recorrer el flujo completo de herramientas, no maximizar desempeño.
\end{definicion}

\begin{definicion}[Bus Avalon Memory-Mapped (Avalon-MM)]
El bus Avalon-MM es el protocolo de interconexión con el que Nios II y sus periféricos se comunican
dentro del sistema. Cada componente es \emph{maestro} (inicia transacciones --- el núcleo Nios II) o
\emph{esclavo} (responde a transacciones en un rango de direcciones --- memoria, JTAG UART, PIO).
Platform Designer arma el interconector (\emph{fabric}) y asigna las direcciones base
automáticamente a partir de los componentes que se agregan al sistema: el diseñador no escribe a
mano la lógica de decodificación de direcciones.
\end{definicion}

\section{El flujo de herramientas: Quartus Prime, Platform Designer, Nios II EDS}

\begin{propiedad}[Tres herramientas, tres responsabilidades]
\begin{center}
\begin{tabular}{|l|p{9.5cm}|}
\hline
\textbf{Herramienta} & \textbf{Responsabilidad} \\ \hline
Quartus Prime & Síntesis, place \& route y programación de la FPGA. Contiene el diseño top-level
  que instancia el sistema Nios II junto con el resto de la lógica del proyecto. \\ \hline
Platform Designer & Ensamblado gráfico del sistema (antes llamado Qsys): agregar el núcleo Nios II
  y sus periféricos, conectar el bus Avalon-MM, asignar direcciones base, y generar el HDL del
  sistema y el archivo \texttt{.sopcinfo} que lo describe. \\ \hline
Nios II EDS & \emph{Eclipse-based Development Suite}: a partir del \texttt{.sopcinfo}, genera el
  BSP (Board Support Package) del sistema exacto que se armó, y compila/descarga/depura software en
  C que corre sobre ese sistema. \\ \hline
\end{tabular}
\end{center}
El orden de dependencia es estricto: Platform Designer necesita el proyecto de Quartus para existir;
Nios II EDS necesita el \texttt{.sopcinfo} que produce Platform Designer --- y ese \texttt{.sopcinfo}
debe corresponder exactamente al bitstream programado en la FPGA, o el software fallará al ejecutar
contra un mapa de memoria que no es el real.
\end{propiedad}

\begin{propiedad}[Nota de versión]
Cyclone IV --- la familia de la DE2-115 --- solo está soportada hasta Quartus Prime 18.1 (Standard
Edition), la última versión que además incluye el flujo clásico de Nios II EDS. Versiones
posteriores de Quartus Prime dejaron de dar soporte tanto a Cyclone IV como al propio Nios II
Classic, en favor de familias más nuevas y de Nios V (basado en RISC-V). Para este laboratorio se
usa Quartus Prime 18.1 Standard Edition.
\end{propiedad}

\section{De hardware a software: BSP y compilación cruzada}

\begin{definicion}[BSP --- Board Support Package]
El BSP es el conjunto de drivers, bibliotecas HAL (\emph{Hardware Abstraction Layer}) y archivos de
configuración que Nios II EDS genera automáticamente a partir del \texttt{.sopcinfo}: expone a nivel
de software exactamente los periféricos, direcciones base e interrupciones del sistema que se armó
en Platform Designer. Si el sistema de hardware cambia --- se agrega un periférico, cambia una
dirección base --- el BSP debe regenerarse; de lo contrario el software queda describiendo un
sistema que ya no coincide con el bitstream programado en la FPGA.
\end{definicion}

\begin{propiedad}[Compilación cruzada]
El programa en C se compila en la máquina host con un compilador cruzado (\texttt{nios2-elf-gcc},
parte de Nios II EDS) que genera código máquina para la arquitectura Nios II, no para la máquina
host. El ejecutable resultante (\texttt{.elf}) se descarga a la memoria del sistema a través del
mismo cable USB-Blaster usado para programar la FPGA, usando el enlace JTAG.
\end{propiedad}

\begin{definicion}[JTAG UART]
La JTAG UART es un periférico Avalon-MM que expone un canal serie sobre el mismo cable JTAG usado
para programar la FPGA --- no necesita un puerto serie físico aparte. Es el mecanismo más simple
para que un primer programa en Nios II se comunique con el computador anfitrión: \texttt{printf} en
C se enruta, vía el HAL del BSP, a la JTAG UART como salida estándar (\texttt{stdout}), y se lee
desde una terminal Nios II en el host.
\end{definicion}

\section{Práctica integradora: sistema Nios II mínimo sobre la DE2-115}

\begin{propiedad}[La tarjeta DE2-115]
La Terasic DE2-115 monta una FPGA Cyclone IV E \texttt{EP4CE115F29C7} (115K elementos lógicos), un
oscilador de 50\,MHz (\texttt{CLOCK\_50}), pulsadores (\texttt{KEY}), interruptores deslizantes
(\texttt{SW}), LEDs rojos y verdes, y un circuito USB-Blaster integrado que sirve tanto para
programar la FPGA como para el enlace JTAG UART --- un único cable USB basta para todo el
laboratorio.
\end{propiedad}

\begin{propiedad}[Componentes del sistema mínimo, en Platform Designer]
\begin{center}
\begin{tabular}{|l|p{9.5cm}|}
\hline
\textbf{Componente} & \textbf{Rol} \\ \hline
Nios II Processor (núcleo \texttt{/e}) & Maestro Avalon-MM; ejecuta el software. \\ \hline
Reloj y reset (exportados) & Entradas de reloj/reset del sistema, conectadas en el top-level de
  Quartus a \texttt{CLOCK\_50} y a un pulsador \texttt{KEY}. \\ \hline
On-Chip Memory (RAM) & Memoria única de instrucciones y datos --- suficiente para un
  \texttt{Hello World}; se le asigna la región de \emph{reset} del núcleo Nios II. \\ \hline
JTAG UART & Canal serie sobre JTAG; destino de \texttt{stdout} del programa en C. \\ \hline
System ID Peripheral & Identificador que Nios II EDS compara entre el \texttt{.sopcinfo} y el
  bitstream programado, para detectar si software y hardware quedaron desincronizados. \\ \hline
\end{tabular}
\end{center}
Platform Designer asigna las direcciones base de cada esclavo automáticamente (menú
\emph{System $\to$ Assign Base Addresses}) al conectar los componentes al bus del núcleo Nios II.
\end{propiedad}

\begin{center}
\begin{tikzpicture}[
  block/.style={draw, rounded corners, minimum width=3.1cm, minimum height=0.9cm, align=center, font=\scriptsize},
  master/.style={block, fill=unicaucaverde!12, minimum width=4.6cm, font=\scriptsize\bfseries},
  bus/.style={draw, unicaucaverde, thick},
]
  \node[master] (nios) at (0,3.3) {Nios II Processor (/e)\\[1pt]\normalfont\scriptsize maestro Avalon-MM};
  \node[block]  (mem)   at (-4.4,0) {On-Chip Memory\\[1pt]\normalfont\scriptsize instrucciones + datos};
  \node[block]  (uart)  at (0,0)    {JTAG UART\\[1pt]\normalfont\scriptsize destino de \texttt{stdout}};
  \node[block]  (sysid) at (4.4,0)  {System ID\\[1pt]\normalfont\scriptsize Peripheral};

  \draw[bus] (-5.1,1.6) -- (5.1,1.6) node[midway, above, font=\scriptsize] {bus Avalon-MM};
  \draw[bus] (nios.south) -- (0,1.6);
  \draw[bus] (mem.north)  -- (-4.4,1.6);
  \draw[bus] (uart.north) -- (0,1.6);
  \draw[bus] (sysid.north)-- (4.4,1.6);

  \draw[-{Latex[length=2mm]}, bus] ($(nios.west)+(-1.3,0.15)$) -- ++(1.3,0)
    node[midway, above, font=\scriptsize] {clk};
  \draw[-{Latex[length=2mm]}, bus] ($(nios.west)+(-1.3,-0.15)$) -- ++(1.3,0)
    node[midway, below, font=\scriptsize] {reset};
\end{tikzpicture}

\textit{\footnotesize Figura 1. Diagrama de bloques del sistema mínimo armado en Platform Designer:
el núcleo Nios II es el único maestro Avalon-MM; memoria, JTAG UART y System ID son esclavos
colgados del mismo bus. Reloj y reset entran al núcleo desde fuera del sistema Qsys --- se exportan
y se conectan en el top-level de Quartus a \texttt{CLOCK\_50} y a \texttt{KEY[0]}.}
\end{center}

\begin{ejemplo}[Procedimiento completo]
\begin{enumerate}[label=\textbf{(\arabic*)}, leftmargin=1.2cm]
  \item \textbf{Proyecto Quartus.} \emph{File $\to$ New Project Wizard}; dispositivo
        \texttt{EP4CE115F29C7} (familia Cyclone IV E), family/device exactos de la DE2-115.
  \item \textbf{Ensamblar el sistema.} \emph{Tools $\to$ Platform Designer}; agregar desde el
        catálogo de IP el núcleo Nios II Processor, On-Chip Memory, JTAG UART y System ID
        Peripheral; conectar las líneas Avalon-MM (clock, reset, datos) entre el núcleo y cada
        esclavo arrastrando en la vista de conexiones; exportar los puertos de reloj y reset del
        sistema hacia el top-level; \emph{System $\to$ Assign Base Addresses}; \textbf{Generate
        HDL}. Esto produce el módulo del sistema (\texttt{.qip}/\texttt{.v}) y el archivo
        \texttt{.sopcinfo} que Nios II EDS necesitará después.
  \item \textbf{Integrar en el top-level.} Instanciar el módulo generado en el archivo top-level de
        Quartus, conectando sus puertos de reloj y reset exportados a las señales físicas de la
        tarjeta.
\begin{lstlisting}[style=vstyle,caption={Instanciación del sistema generado en el top-level (Verilog)}]
module de2_115_niosii_top (
    input  wire CLOCK_50,
    input  wire [3:0] KEY
);

    niosII_system u0 (
        .clk_clk        (CLOCK_50),
        .reset_reset_n  (KEY[0])   // KEY es activo en bajo: presionado = reset
    );

endmodule
\end{lstlisting}
        El nombre exacto del módulo y de sus puertos (\texttt{clk\_clk}, \texttt{reset\_reset\_n})
        depende de cómo se hayan nombrado las interfaces exportadas en Platform Designer; Quartus
        muestra el nombre real generado en el panel \emph{IP Catalog $\to$ Instantiation Template}.
  \item \textbf{Asignar pines.} Usar la asignación de pines \emph{oficial} de Terasic para la
        DE2-115 (archivo \texttt{.csv}/\texttt{.qsf} de asignación de pines que acompaña el manual
        de la tarjeta), importándola con \emph{Assignments $\to$ Import Assignments} en vez de
        escribir los números de pin a mano --- así se evita el error más común y más difícil de
        depurar en un primer laboratorio de FPGA: un pin mal escrito que dejar la tarjeta
        ``sin responder'' sin ningún mensaje de error de Quartus.
  \item \textbf{Compilar y programar.} \emph{Processing $\to$ Start Compilation} (Analysis \&
        Synthesis, Fitter, Assembler); luego \emph{Tools $\to$ Programmer}, seleccionar el cable
        USB-Blaster y el archivo \texttt{.sof} generado, y programar la FPGA.
  \item \textbf{Software en Nios II EDS.} Abrir la \emph{Nios II 18.1 Command Shell} o el IDE
        Eclipse de Nios II EDS; \emph{File $\to$ New $\to$ Nios II Application and BSP from
        Template}, apuntando al \texttt{.sopcinfo} generado en el paso (2) y eligiendo la plantilla
        \texttt{Hello World}.
\begin{lstlisting}[caption={\texttt{hello\_world.c}, tal como lo genera la plantilla}]
#include <stdio.h>

int main(void) {
    printf("Hello from Nios II!\n");
    return 0;
}
\end{lstlisting}
  \item \textbf{Compilar y ejecutar.} \emph{Build Project} (compila BSP y aplicación); \emph{Run As
        $\to$ Nios II Hardware} descarga el \texttt{.elf} a la memoria on-chip por JTAG y lo
        ejecuta; el mensaje aparece en la consola Nios II del IDE (equivalente a correr
        \texttt{nios2-terminal} desde la Nios II Command Shell).
\end{enumerate}
\end{ejemplo}

\begin{propiedad}[Fallas comunes y cómo leerlas]
\begin{itemize}[leftmargin=1.2cm]
  \item \textbf{No hay salida en la terminal Nios II}: casi siempre el reloj o el reset del sistema
        no quedaron conectados en el top-level, o se conectaron con la polaridad invertida (recordar
        que \texttt{KEY} en la DE2-115 es activo en bajo).
  \item \textbf{Error de ``System ID mismatch'' al ejecutar}: el \texttt{.sof} programado en la FPGA
        no corresponde al \texttt{.sopcinfo} usado para generar el BSP --- se modificó el sistema en
        Platform Designer sin recompilar y reprogramar el proyecto de Quartus. Se corrige
        recompilando y reprogramando antes de volver a ejecutar el software.
  \item \textbf{Tiempo de espera agotado al programar}: el cable USB-Blaster no está bien conectado,
        o la tarjeta no tiene alimentación --- revisar antes de sospechar del proyecto.
\end{itemize}
\end{propiedad}

\section{Ejercicios propuestos}

\begin{enumerate}[leftmargin=1.2cm]
  \item Explica, en tus palabras, por qué el archivo \texttt{.sopcinfo} debe corresponder
        exactamente al \texttt{.sof} programado en la FPGA, y qué error concreto aparece cuando no
        coinciden.
  \item ¿Por qué el núcleo Nios II/e --- el más simple, sin caches --- es una elección razonable
        para el sistema de esta semana, y en qué situación futura convendría cambiar a Nios II/f?
  \item Agrega un componente PIO de salida (\emph{Parallel I/O}, ancho 1 bit) al sistema en Platform
        Designer, conéctalo al bus, regenera el sistema y el BSP, y modifica
        \texttt{hello\_world.c} para que además de imprimir el mensaje, encienda un LED de la
        DE2-115 escribiendo sobre el registro de datos del PIO.
\end{enumerate}

\end{document}
